\documentclass[compress]{beamer}
\mode<presentation>
{
 \usetheme{Vilanova}
}

\usepackage[french]{babel}

\usepackage[utf8]{inputenc}

\usepackage{times}
\usepackage[T1]{fontenc}

\usepackage{amsfonts}
\usepackage{amsmath}
\usepackage{amssymb}
\usepackage{tikz}
\usepackage{eurosym}
%\usepackage{url}
\usepackage[normal]{subfigure}
\newcommand{\goodgap}{%
	\hspace{\subfigtopskip}%
	\hspace{\subfigbottomskip}}



%\newtheorem{definition}{Definition}

\title{Traitement d'images de télédétection}
\subtitle{Corrections géométriques}



\author
{jordi.inglada@cesbio.cnes.fr}
\normalsize

\institute[Cesbio] % (optional, but mostly needed)
{\textsc{Centre d'Études Spatiales de la Biosphère, Toulouse, France}}

\date{}

\pgfdeclareimage[height=96mm,width=128mm]{background}{fondsClairSansLogo}
\setbeamertemplate{background}{\pgfuseimage{background}}
\pgfdeclareimage[height=0.6cm]{logoIncrust}{logoIncrust}
\pgfdeclareimage[height=0.5cm]{logo_cesbio}{logo_cesbio}
\logo{
\begin{tabular}{lp{0.25\textwidth}lp{0.25\textwidth}r}
\href{http://www.cesbio.ups-tlse.fr/}{\pgfuseimage{logo_cesbio}}
&&\footnotesize{AUF - Marrakech 2011}&&
\href{http://www.orfeo-toolbox.org}{\pgfuseimage{logoIncrust}}\\
\end{tabular}
}


\subject{Image geometry in ORFEO Toolbox}




% Delete this, if you do not want the table of contents to pop up at
% the beginning of each subsection:
\AtBeginSubsection[]
{
  \begin{frame}<beamer>
    \frametitle{Outline}
    \tableofcontents[currentsection,currentsubsection]
  \end{frame}
}




% If you wish to uncover everything in a step-wise fashion, uncomment
% the following command: 

%\beamerdefaultoverlayspecification{<+->}
\begin{document}

\begin{frame}
  \titlepage
  \begin{center}
{\tiny Ce contenu est dérivé de la formation \href{http://www.orfeo-toolbox.org/packages/PragmaticRemoteSensing-handout.pdf}{``Pragmatic Remote
  Sensing''} dispensée par J. Inglada et E. Christophe en juillet 2010
  dans le cadre du colloque IGARSS. Il est mis à disposition selon les termes de la licence :\\
Creative Commons Paternité – Partage à l’Identique 3.0 non transcrit.} \href{http://creativecommons.org/licenses/by-sa/3.0/}{\includegraphics[width=0.05\textwidth]{/home/inglada/Dev/GH/IGARSS2010/Tutorial/Slides/Ressources/CC-licence.png}}    
  \end{center}
\end{frame}


\section*{Introduction}

\begin{frame}

  \frametitle{Introduction}
  \hspace*{-1cm}
  \begin{tikzpicture}[scale=0.165]
    \tiny
    \draw[fill=black!10] (-1,-12) rectangle (75,17);
     \foreach \x in {5,...,1}
       \draw[fill=red] (\x,\x) rectangle +(4,4);
     \node[fill=black!10, text width= 1.5cm] (InputSeries) at
       (4,-1) {Série d'images};
     \pause
     \draw[->,thick] (9,5) --  +(3,0);
     \pause
     \draw[fill=black!30,rounded corners=2pt] (12.2,3) rectangle +(6,4);
     \node[text width= 0.8cm] (SensorModel) at (15,5) {Modèle capteur};
     \pause
     \draw[fill=red!30] (1,-10) rectangle +(4,4);
     \node[fill=black!10, text width= 1.2cm] (DEM) at
       (5,-11) {MNT};
     \pause
     \draw[->,thick] (3,-5.5) --  ++(0,3) -- ++(12,0) -- ++(0,5);
     \pause
     \draw[->,thick] (18.5,5) --  +(3,0);
     \pause
     \foreach \x in {5,...,1}
       \draw[fill=blue,xshift=600pt] (\x,\x) rectangle +(4,4);
     \node[fill=black!10, text width= 2.8cm] (GeoRefSeries) at
       (28,-1) {Géo-référencement};
\pause
      

       \draw[->,thick] (25.5,8.5) --  +(0,3);
       
     \draw[fill=black!30,rounded corners=2pt] (22,12) rectangle +(8.5,4);
     \node[text width= 1.5cm] (HomPoExtr) at (27,14) {Points Homologues};

     \draw[->,thick] (21.5,14) --  +(-2.5,0);

     \draw[fill=black!30,rounded corners=2pt] (11,12) rectangle +(8,4);
     \node[text width= 1.3cm] (BBAdj) at (15.5,14) {Spatio triangulation};

     \draw[->,thick] (15,11.5) --  +(0,-4);

     \pause
      \draw[->,thick] (30,5) --  +(3,0);
      \pause
     \draw[fill=black!30,rounded corners=2pt] (33.2,2.5) rectangle +(6,4.5);
     \node[text width= 0.7cm] (FineRegistration) at (36,4.9) {Recalage fin};
     \pause

     
     \draw[->,thick] (39.5,5) --  +(3,0);
     \pause
     \foreach \x in {5,...,1}
       \draw[fill=green,xshift=1200pt] (\x,\x) rectangle +(4,4);
     \node[fill=black!10, text width= 1.8cm] (RegistSeries) at
       (47,-1) {Série recalée};
     \pause
     \draw[->,thick] (36,2) --  ++(0,-10) -- ++(-30,0);

     \pause
      \draw[->,thick] (52,5) --  +(3,0);
      \pause
     \draw[fill=black!30,rounded corners=2pt] (55.2,2.5) rectangle +(6,4.5);
     \node[text width= 0.7cm] (CartoProjection) at (57.5,4.9)
          {Projection Carto};
     \pause

     
     \draw[->,thick] (61.5,5) --  +(3,0);
     \pause
     \foreach \x in {5,...,1}
       \draw[fill=yellow,xshift=1810pt] (\x,\x) rectangle +(4,4);
     \node[fill=black!10, text width= 1.95cm] (CartoSeries) at
       (68,-1) {Ortho-images};
     
       
  \end{tikzpicture}
\end{frame}

%% \begin{frame}
%%   \frametitle{Introduction}
%%  \begin{block}{How to register image series}
%%  \begin{itemize}
%%  \item Sensor models and bundle-block adjustment
%%  \item Homologous point extraction
%%  \item Fine registration
%%  \end{itemize}
%%  \end{block}
%%  \begin{block}{How to measure the quality}
%%  \begin{itemize}
%%  \item Choosing the reference
%%  \item Quality measures
%%  \end{itemize}
%%  \end{block}
%%  \end{frame}


\section[Modèles]{Modèles de capteur}


\begin{frame}
  \frametitle{Modèles de capteur}

  \framesubtitle{Définition}
Transformation de coordonnées entre l'image issue du capteur $(l,c)$
et les coordonnées au sol $(X,Y)$ pour chaque pixel :
\pause
\begin{displaymath}
  \begin{array}{cc}
    Direct & \\
    X = f_x(l,c,h,\vec\theta) & Y = f_y(l,c,h,\vec\theta)\\
     & \\ \pause
    Inverse & \\
    l = g_l(X,Y,h,\vec\theta) & c = g_c(X,Y,h,\vec\theta)
  \end{array}
\end{displaymath}
\pause
Où $\vec\theta$ est l'ensemble de paramètres décrivant le capteur et
la géométrie d'acquisition.\\
\pause
L'élévation de chaque point (MNT) doit être connue.
  
\end{frame}

\begin{frame}
  \frametitle{Modèles de capteur}

  \framesubtitle{Types de modèles}
  \begin{itemize}
    \item Modèles physiques
      \begin{itemize}
	\item Rigoureux, complexes, équations fortement non-linéaires
	\item Difficiles à inverser
	\item Les paramètres ont une signification physique
	\item Spécifiques à chaque capteur
      \end{itemize}
    \item Modèles analytiques génériques
      \begin{itemize}
	\item Ex: polynomiaux, fractions rationnelles, etc.
	\item Moins précis
	\item Faciles à mettre en oeuvre
	\item Les paramètres peuvent ne pas avoir de signification physique
      \end{itemize}
  \end{itemize}

\end{frame}

\begin{frame}
  \frametitle{Modèles de capteur}

  \framesubtitle{L'approche OTB}
  \begin{itemize}
    \item Utilisation de {\em factories} : les modèles sont générés
      automatiquement en utilisant les méta-données des images
    \item Modèles disponibles
      \begin{itemize}
	\item Fractions rationnelles : Quickbird, Ikonos, WorldView-2
	\item Modèles physiques : SPOT5
	\item Radar : ERS, ASAR, Radarsat, Cosmo Skymed, TerraSAR-X, Palsar
      \end{itemize}
  \end{itemize}
\end{frame}


\begin{frame}
\frametitle{La main à la pâte}
\begin{enumerate}
\item Monteverdi : Ouvrir une image en géométrie capteur
  \begin{itemize}
  \item /home/auf/OTB/data/TOULOUSE (image Quickbird de Toulouse)
  \item Choisir le fichier TIF dans l'un des 2 répertoires
  \end{itemize}
\item Afficher l'image
  \begin{itemize}
  \item Click droit sur le nom de l'image dans la fenêtre principale
    de Monteverdi
  \item ``Display in viewer''
  \item Attendre la génération du quicklook
  \end{itemize}
\item Observer comment les coordonnées géographiques sont recalculées
  quand le curseur se déplace
  \begin{itemize}
  \item Elles sont mises à jour en appliquant le modèle de capteur
    crée à partir des méta-données de l'image
  \end{itemize}
\end{enumerate}
\end{frame}

\begin{frame}
  \frametitle{Modèles de capteur}

  \framesubtitle{Utilisation : ortho-rectification}
  \begin{enumerate}
    \item Lecture des méta-données image et création du modèle avec
      les bons paramètres
  \item Définition de la ROI en coordonnées sol (c'est la matrice de
    pixels de sortie)
  \item Balayer les pixels de coordonnées $(X,Y)$ :
    \begin{enumerate}
      \item Obtenir $h$ à partir du MNT
      \item Calculer $(c,l) = G(X,Y,h,\vec\theta)$
      \item Interpoler les valeurs des pixels si $(c,l)$ ne sont pas
        des valeurs entières
    \end{enumerate}
  \end{enumerate}
\end{frame}

\begin{frame}
\frametitle{La main à la pâte}
\begin{enumerate}
\item Monteverdi: Geometry $\rightarrow$ Reproject image
\item Choisir l'image à ortho-rectifier (Panchromatique)
\item Choisir les paramètres
  \begin{itemize}
  \item Onglet ``Output image''
    \begin{itemize}
    \item Par défaut, projection UTM.
    \item Choisir une petite taille (size x, size y) pour que ça aille
    vite!
    \end{itemize}
  \item Regarder l'onglet ``Input image'' : l'image en entrée est bien
    en géométrie capteur.
  \end{itemize}
\item Sauvegarder le résultat

\end{enumerate}
\end{frame}


\begin{frame}
  \frametitle{Modèles de capteur}
  \framesubtitle{Limites de l'approche}

  \begin{itemize}
    \item Un géo-référencement précis nécessite :
      \begin{itemize}
	\item Un MNT précis
	\item Des paramètres capteur sans erreur, $\vec\theta$
      \end{itemize}
    \item Pour les séries multi-temporelles d'images on a besoin de
      \alert{recalage fin} :
      \begin{itemize}
      \item Précision sous-pixellique
      \item Pour chaque pixel de la scène
      \end{itemize}
    \item Les MNT et les méta-données capteur ne fournissent pas cette
      précision.
    \item Solution : utilisation de l'information redondante entre les
      images de la série.
  \end{itemize}
\end{frame}

\section{Optimisations}
\subsection{Spatio-triangulation}

\begin{frame}
  \frametitle{Spatio-triangulation}
  \framesubtitle{Position du problème}
  \begin{columns}[T]
\column{.5\textwidth}
  \begin{itemize}
    \item La série d'images est ortho-rectifiée (avec le MNT et les
      paramètres disponibles).
    \item Supposons que des point homologues (PH) peuvent être obtenus
      aisément : $PH_i = (X_i,Y_i,h_i)$
    \item Pour chaque image et pour chaque point nous pouvons écrire :
    $(l_{ij},c_{ij}) = G_j(X_i,Y_i,h_i,\vec\theta_j)$
  \end{itemize}
\column{.5\textwidth}
\begin{tikzpicture}[scale=0.15]
\draw[fill=yellow!20] (-5.5,-15.5) rectangle (5.5,-5.5);
    \draw[step=0.5, gray, very thin] (-5.5,-15.5) grid (5.5,-5.5);

    \draw[fill=green!20,rotate=10] (-15.5,0.5) rectangle (-5.5,10.5);
    \draw[step=0.5, gray, very thin,rotate=10] (-15.5,0.5) grid
    (-5.5,10.5);

    \draw[fill=blue!20,rotate=-10] (5.5,0.5) rectangle (15.5,10.5);
    \draw[step=0.5, gray, very thin,rotate=-10] (5.5,0.5) grid
    (15.5,10.5);

    \pause
    \draw[fill=red!70] (1,-11) circle (0.2);
    \pause
    \draw (1,-11) .. controls +(30:1cm) and +(60:1cm) .. (-10,7);
    \pause
    \draw[fill=red!70] (-10,7) circle (0.2);
    \pause
    \node (eq1) at (-12.2,-4) {$\scriptstyle{G_1(X_i,Y_i,h_i,\vec\theta_1)}$};
    \pause
    \draw (1,-11) .. controls +(-30:1cm) and +(-60:1cm) .. (10,7);
    \pause
    \draw[fill=red!70] (10,7) circle (0.2);
    \pause
    \node (eq2) at (7.2,-3) {$\scriptstyle{G_2(X_i,Y_i,h_i,\vec\theta_2)}$};
    
\end{tikzpicture}
\begin{itemize}
      \item Tout est connu.
\end{itemize}
  \end{columns}
\end{frame}

\begin{frame}
  \frametitle{Spatio-triangulation}
  \framesubtitle{Affinage du modèle}
  \begin{itemize}
    \item Si nous définissons $\vec\theta_j^R = \vec\theta_j +
    \vec{\Delta\theta_j}$ comme étant les paramètres affinés,
    $\vec{\Delta\theta_j}$ ce sont les inconnues du problème d'affinage.
    \item Nous avons beaucoup plus d'équations que d'inconnues si nous
      disposons de beaucoup de PH.
    \item Solution par moindres carrés
      \begin{itemize}
	\item Nous avons besoin des dérivées du modèle de capteur par
          rapport à ses paramètres.
      \end{itemize}
  \end{itemize}
  
\end{frame}

\begin{frame}
  \frametitle{La main à la pâte}
  \framesubtitle{Recalage manuel de 2 images}
\vspace*{-0.6cm}
{\scriptsize
  \begin{itemize}
  \item Monteverdi : Geometry $\rightarrow$ Homologous points extraction
  \item Choisir 2 images avec une zone commune
    \begin{itemize}
    \item /home/auf/OTB/data/Examples/QB\_Suburb.png et QB\_SuburbR10X13Y17.png
    \end{itemize}
  \item L'IHM permet de choisir la transformation géométrique
    \begin{itemize}
    \item Choisir Affine
    \end{itemize}
  \item On peut sélectionner des PH dans la zone de zoom et les
    ajouter à la liste
  \item Quand on a choisi plusieurs PH, on peut évaluer la transformation
  \item On peut ensuite utiliser le bouton {\em guess} afin de prédire
    la position des nouveaux points
  \item L'IHM affiche les paramètres de la transformation estimée,
    l'erreur commise sur chaque point et l'EQM
  \item On peut éliminer de la liste les points qui ont le plus
    d'erreur
  \end{itemize}
  }
\end{frame}





\begin{frame}
  \frametitle{La main à la pâte}
  \framesubtitle{Projection sur une autre image}
  \begin{itemize}
  \item Monteverdi: Geometry $\rightarrow$ Superimpose 2 Images
  \item Choisir une image quelconque comme image à reprojeter
  \item Choisir une image ortho-rectifiée comme référence
    \begin{itemize}
    \item S'assurer que les 2 images ont une zone commune!
    \end{itemize}
  \item Utiliser le même MNT que celui utilisé pour l'image ortho-rectifiée
  \item Save/Quit
  \end{itemize}
\end{frame}


\end{document}
